\chapter{数值计算考量}

\begin{solutionexercise}{1}
\problemheading 为一种 6 位格式（1 位符号、3 位尾数、2 位指数）绘制与原书图 A.5
等价的图。用结果说明尾数每增加一位，会如何改变数轴上的可表示数集合。作为必要上下文，
图 A.5 的基准是 1 位符号、2 位显式尾数和 2 位指数的教学“无零”格式：指数偏置为 1，
指数编码 11 保留，不采用零或反规格化数，只绘正数，负数关于 0 对称。新图须列出全部有效
指数区间内的点及相邻间距。

\approachheading 沿用原书图 A.5 的教学格式：指数 11 保留，指数偏置为 1，不采用
反规格化数；只列正数，负数关于 0 对称。三位显式尾数 $M=k/8$，$k=0,\ldots,7$。

\answerheading 对有效指数编码 00、01、10，实际指数分别为 $-1,0,1$。正可表示数为
\[
\begin{array}{c|l|c}
E&\text{正可表示数}&\text{相邻间距}\\ \hline
-1&0.5,0.5625,0.625,0.6875,0.75,0.8125,0.875,0.9375&1/16\\
0&1,1.125,1.25,1.375,1.5,1.625,1.75,1.875&1/8\\
1&2,2.25,2.5,2.75,3,3.25,3.5,3.75&1/4
\end{array}
\]
负集合是这些数取负；该“无零”格式仍不能表示 0。一般地，每个指数区间的值是
\[
  (1+k/2^m)2^E,\qquad k=0,\ldots,2^m-1.
\]
尾数从 2 位增至 3 位后，每个区间的点数由 4 翻倍到 8，相邻间距和 round-to-nearest
在该区间内的最大绝对误差都减半；指数区间和数量级范围没有改变。换言之，增加尾数位提高
精度，不负责扩大指数范围。

\conclusionheading 与原书图 A.5 相比，在每一对原有相邻点的中点插入一个新点；正负两侧
均如此。
\discussionheading
\discussionpoint{玩具格式首先是用于隔离概念的模型。}
为什么不直接把这 6 位当作一个缩小的 IEEE 格式？本题刻意采用显式尾数，并删除零、反规格化数、无穷和 NaN，目的在于让“增加一位尾数”成为唯一变量，因而可以逐个位型枚举。IEEE binary32 则使用隐含首位、偏置指数和保留编码；套用那套解码规则会同时改变端点、点数与原点附近的结构，所得数轴已不是题目要求的实验。这样的简化模型很适合解释因果，但引用其结论时必须同时声明省略项；在数值格式设计中，先隔离变量、再恢复特殊值和异常路径，也是一种常用的分层验证方法。

\discussionpoint{尾数位决定一个 binade 内的采样密度。}
如果指数保持为 $E$，多出的一位究竟改变了哪里？对含 $p$ 个显式小数位的本题格式，区间 $[2^E,2^{E+1})$ 内的相邻间距为 $2^{E-p}$，所以 $p$ 每增加 1，点数翻倍、间距和最近舍入的最大局部绝对误差都减半。指数端点并未移动；当 $E$ 增加 1 时，绝对 ULP 又随尺度翻倍，而正规格化数的相对间距仍约为 $2^{-p}$。这一区分解释了浮点数为何能在很宽的数值范围内维持近似恒定的相对精度，也提醒工程人员不要用某个固定绝对容差评判所有数量级的数据。

\discussionpoint{零附近的表示揭示渐进下溢的作用。}
若计算结果从最小正规格化数继续趋近于零，本题格式会怎样？因为它既没有零也没有反规格化数，数轴在原点两侧留下空洞；IEEE 的反规格化数则固定指数并逐步减少有效首位，用丢失相对精度的代价保持固定的最小绝对间距，使结果能够渐进地下溢到零。若硬件采用 flush-to-zero，同一个微小差值会更早变成零，迭代停止条件、梯度传播和守恒量都可能随之改变。因此比较低精度格式时不能只列最大值，还应分别检查最小正规格化数、最小非零值、反规格化支持、零的符号及异常值语义。

\discussionpoint{低精度设计是在有限位预算内分配风险。}
同为 8 位或 6 位，为什么不同应用会选择不同的指数与尾数比例？训练梯度可能跨越许多数量级，更需要指数范围；量化后的权重若已按通道缩放，主体误差往往更受尾数密度支配，而块浮点还让一组数共享指数以节省元数据。输入格式也不能单独决定结果质量：乘积常以低精度产生，却在 FP16 或 FP32 中累加，以免长归约持续丢失低位。实际选型因而要结合张量分布、缩放粒度、饱和与非有限值计数、累加器精度和端到端任务指标；“再加一位尾数”只有在瓶颈确由量化噪声主导时才值得付出范围或存储成本。
\variantheading 若尾数再增加到 4 位，则每个指数区间有 16 个正数；$E=0$ 时依次为
$1,1+1/16,\ldots,1+15/16$，间距为 $1/16$，round-to-nearest 的最大绝对误差为
$1/32$。指数范围仍是 $-1,0,1$。
\practiceheading 设计自定义低精度格式时，应写一个穷举编码/解码测试，逐个位型验证单调性、
相邻间距、舍入边界和饱和行为。工程上还要给累加器单独选精度；低精度输入并不意味着应当
用同样低的精度累加。
\sourcecheck{https://standards.ieee.org/ieee/754/6210/}{IEEE 754-2019 官方标准页面；对抗检查：图 A.5 使用的是刻意简化的“无零”格式，不能把 IEEE 的零、反规格化数、无穷和 NaN 位模式混入本题枚举}
\end{solutionexercise}

\begin{solutionexercise}{2}
\problemheading 为另一种 6 位格式（1 位符号、2 位尾数、3 位指数）绘制与原书图 A.5
等价的图。用结果说明指数每增加一位，会如何改变数轴上的可表示数集合。沿用同一教学
“无零”约定：三位指数的偏置为 3，编码 111 保留，不采用零或反规格化数，只绘正数，负数
关于 0 对称；须列出全部有效指数区间，并区分动态范围变化和每个区间内的相对精度。

\approachheading 三位指数偏置为 3，编码 111 保留，所以实际指数为 $-3$ 到 $3$；两位
尾数给每个指数区间 4 个点。

\answerheading 正可表示数为
\[
\begin{array}{c|l}
E&\text{正可表示数}\\ \hline
-3&0.125,0.15625,0.1875,0.21875\\
-2&0.25,0.3125,0.375,0.4375\\
-1&0.5,0.625,0.75,0.875\\
0&1,1.25,1.5,1.75\\
1&2,2.5,3,3.5\\
2&4,5,6,7\\
3&8,10,12,14
\end{array}
\]
负数仍关于 0 对称，0 仍缺失。指数从 2 位增至 3 位后，有效指数层数从
$2^2-1=3$ 增至 $2^3-1=7$。相对于原书图 A.5 的 $E=-1,0,1$，新增
$E=-3,-2,2,3$ 四个区间：最小正规格化数从 $2^{-1}$ 降到 $2^{-3}$，缩小 4 倍；
最大正数从 $1.75\,2^1=3.5$ 扩到 $1.75\,2^3=14$，增大 4 倍。每个区间仍只有 4 个点，
所以相对精度不变。

\conclusionheading 指数位主要扩展动态范围；它不会增加同一二进制数量级内的采样密度。
\discussionheading
\discussionpoint{指数位以倍增编码换取指数级值域。}
给指数增加一位，为什么最大值不是简单地乘 2？指数有 $k$ 位时至多提供 $2^k$ 个编码；新增一位近似把可选指数层数翻倍，而每相邻一层又把数值尺度乘 2，于是端点按新增的指数跨度呈指数扩张。本题从三个有效层扩为七个层，最大指数由 1 到 3，故最大有限值乘 $2^{3-1}=4$；最小端点也按相反方向缩小 4 倍。偏置只决定编码区间放在数轴何处，保留编码、特殊值和反规格化规则才共同决定真实范围，这也是阅读任何硬件格式说明时必须逐项核对的边界。因此端点变化应从实际指数集合推导，而非凭位数口号估计。

\discussionpoint{固定位宽迫使范围与精度此消彼长。}
若芯片只能为每个数保留固定的总位数，多一位指数从哪里来？除符号位外，它通常要占用原本属于尾数的一位，使溢出和下溢变少，却把每个 binade 内的点数减半、相对舍入误差界加倍；本题暂时保持尾数位数不变，只是为了单独观察指数效果。具有少量离群值的数据可能需要宽范围，但若为这些离群值牺牲全部主体精度，整体均方误差反而会上升。格式选择因此是风险分配问题：既要统计极端值造成的饱和损失，也要估计高概率区间的量化噪声，并考虑误差经过后续归约或非线性算子后会怎样放大。

\discussionpoint{缩放让软件参与数值格式设计。}
硬件格式的指数不够时，软件是否只能接受溢出？损失缩放会在反向传播前放大梯度、结束后再除回去，张量或通道量化则保存 scale，把业务值域映射到编码值域；二者都等价于为一批数补充一个共享指数。设量化上界为 $q_{\max}$、观测绝对最大值为 $a$，常见初始尺度约取 $a/q_{\max}$，但单个异常值会让大多数编码挤在零附近。动态统计、异常值旁路和更细缩放粒度能缓解此问题，却增加元数据、同步与复现成本；因此工程上应把缩放策略视为数值算法的一部分，而非格式选定后的无关预处理。

\discussionpoint{FP8 变体说明格式必须匹配数据角色。}
为什么同一套模型可能同时使用不止一种 8 位浮点格式？范围宽、尖峰多的梯度更容易受溢出支配，倾向保留更多指数位；经过归一化或离线标定的权重与激活范围较窄，可能更需要尾数位来降低局部量化误差。即使输入都为 FP8，乘加的部分和通常仍提升到 FP16 或 FP32，因为数千项归约的舍入行为不同于一次乘法。部署验证应按张量角色记录饱和、下溢、非有限值及任务损失，并用真实层分布而非均匀随机数做压力测试；这种按数据角色分配格式的思想也延伸到混合精度求解器和分层存储系统。
\variantheading 若指数扩为 4 位、偏置取 7、编码 1111 保留，同时仍用两位尾数，则有效
指数为 $-7$ 到 $7$。最小正数为 $2^{-7}=1/128$，最大正数为
$1.75\times2^7=224$；每个 binade 仍只有 4 个点。
\practiceheading 混合精度训练和推理常用缩放因子把张量值域移入有限指数范围。上线前应记录
溢出、下溢和饱和计数，并把缩放策略与模型版本绑定；仅看平均误差可能掩盖少量但致命的
溢出值。
\sourcecheck{https://standards.ieee.org/ieee/754/6210/}{IEEE 754-2019 官方标准页面；对抗检查：不能简单说“范围翻倍”，新增一位指数使指数编码数翻倍，并使数值端点按 2 的新增指数跨度成倍扩展}
\end{solutionexercise}

\begin{solutionexercise}{3}
\problemheading 假设一种新处理器受技术困难限制，执行加法的浮点算术单元只能“向零
舍入”（通过向 0 截断数值来舍入）。硬件保留了足够多的位，唯一误差来自最终舍入。这台
机器上加法操作的最大 ULP 误差是多少？说明“最大值”是可达到值还是上确界，并指出溢出
等不属于题设模型的边界。

\approachheading 在同一 binade 内取两个相邻可表示数 $a<b$，间距定义为 1 ULP。向零
舍入总选择绝对值较小的一侧，而不是最近的一侧。

\answerheading 对正的精确结果 $x\in[a,b)$，结果为 $a$，故
\[
  0\le x-\operatorname{fl}_{0}(x)<b-a=1\ \mathrm{ULP}.
\]
负数情形对称。因而严格说误差的\textbf{上确界}是 1 ULP；任一有限、正常且在表示范围内
的非精确结果误差都严格小于 1 ULP，但可以任意逼近它。教材或硬件规格把上界写成“最大
1 ULP”时，指的就是这个上确界。相比之下，round-to-nearest 的对应上界是 0.5 ULP。

这个结论排除溢出；若精确和超出有限表示范围，错误不能再由局部 1 ULP 界描述。在最小
反规格化数区域，ULP 取固定的最小间距，同样适用。

\conclusionheading 向零舍入的误差界为 $<1$ ULP，通常报告为最大 1 ULP，而不是
0.5 ULP。
\discussionheading
\discussionpoint{ULP 是随数轴位置伸缩的局部刻度。}
“误差 1 ULP”能否换算成一个固定的小数？不能，因为 ULP 通常取结果附近两个相邻浮点数的间距：在同一 binade 内它保持为 $2^{E-p}$，越过二的幂边界后便翻倍；反规格化区则固定为最小间距。于是绝对误差同为 $10^{-6}$，在不同数量级上可能对应完全不同的 ULP 数，而相同的 1 ULP 也可能有不同绝对大小。ULP 特别适合检验实现是否返回附近的可表示数，但报告时必须说明参考点、格式及特殊区域；面向应用的误差预算还应同时给出绝对误差、相对误差或问题特定范数。

\discussionpoint{舍入模式同时决定误差大小和方向。}
两个相邻浮点数之间的精确结果应落向哪一端？向零舍入总选绝对值较小的一端，所以正数误差向下、负数误差向上，任一有限非精确结果的误差严格小于 1 ULP，但可无限逼近该界；向正无穷和负无穷则分别给出定向包围。区间算法会把同一表达式分别在向负无穷和向正无穷模式下求值，前者提供不大于精确值的下界，后者提供不小于精确值的上界，从而让真实结果即使经过多步运算仍被括住。最近偶数在非溢出区域把界缩为 0.5 ULP，并在恰逢中点时按末位奇偶打破平局，以避免总朝同一方向累积偏差。定向舍入因此适合区间算术和保守几何边界，而统计计算常偏好近似无偏的最近舍入，选择模式其实是在准确度、方向保证和硬件成本之间取舍。

\discussionpoint{局部舍入界不能替代问题条件分析。}
若每次加法都只错不到 1 ULP，整个算法是否就一定准确？设两个输入已分别带有微小误差，再计算 $10^8-10^8$ 一类严重抵消，结果的绝对误差可能仍小，但相对于真实的小差值却会非常大；局部正确舍入并不能恢复输入中已经丢失的信息。数值分析因而区分前向误差、后向误差与条件数：前两者描述实现偏离结果或等价输入多少，条件数描述问题本身对扰动有多敏感。工程上需要把硬件操作界沿算法传播，并用残差、缩放或稳定重写控制放大效应，不能把“每步 1 ULP”误读成端到端质量承诺。

\discussionpoint{舍入测试必须精确命中边界。}
随机生成一百万组加数，为什么仍可能漏掉舍入错误？最难的路径往往只出现在两个相邻浮点数的中点附近、binade 交界、反规格化边界或溢出前一格，随机实数几乎不会精确命中这些位型。可靠测试可用整数位模式构造相邻数，以多精度算术形成中点及其两侧的值，并分别覆盖正负号、精确结果、进位和 tie-to-even。还要防止编译器把乘加收缩为 FMA 或重排表达式，从而改变中间舍入点；把编译选项、实际指令和参考舍入模式纳入测试记录，才能区分算术单元错误与工具链改变语义。
\variantheading 若改为向 $+\infty$ 舍入，有限非精确结果的绝对误差上确界仍是 1 ULP，
但误差方向总为非负：正数向绝对值更大处移动，负数向绝对值更小处移动。若改为最近偶数
舍入，上确界降为 0.5 ULP。
\practiceheading 定向舍入常用于区间算术构造保守上下界。实现时必须确认编译器没有重排或
融合破坏舍入模式假设；CUDA 代码可使用带明确舍入后缀的算术内建函数，并用相邻浮点数测试
每个边界。
\sourcecheck{https://docs.nvidia.com/cuda/floating-point/index.html}{NVIDIA CUDA Floating Point 官方资料对 IEEE 舍入模式和 ULP 的说明；对抗检查：把“1 ULP”说成每个非精确结果都能恰好达到，或忽略溢出，均不严谨}
\end{solutionexercise}

\begin{solutionexercise}{4}
\problemheading 一名研究生编写了 CUDA 内核，把一个大型浮点数组归约为所有元素之和。
数组总是按数值从小到大排序。为避免分支分歧，他决定实现原书图 10.7 的算法。该低分歧
分配在首轮令前半数组的线程 $i$ 计算 $x_i+x_{i+N/2}$，以后继续折半。请解释这为何可能
降低结果的准确度，并把浮点加法次序、ULP 和结合律纳入说明。

\approachheading 原书图 10.7 首轮把前半数组与后半数组逐项配对，而不是让相邻元素配对。
对已排序的非负数据，这会过早地把小数加到大数上。

\answerheading 设数组长 $2n$。该分配首轮计算
\[
  x_0+x_n,\ x_1+x_{n+1},\ldots,x_{n-1}+x_{2n-1}.
\]
若 $x_i$ 与 $x_{n+i}$ 的指数差很大，较小加数在对齐尾数后可能完全落到保留位之外；当它
小于大数处的半个 ULP 时，round-to-nearest 甚至直接返回大数。大量小贡献会在尚未彼此
聚合时分别丢失。若先把数量级接近的相邻小数做平衡树归约，小数先形成较大的部分和，之后
再与大数相加，通常能保留更多有效位。补偿求和或按绝对值分桶还能进一步改善。

这里的“降低”是可能性而不是逐个输入的定理：若数值同量级、加法恰好可表示，或正负抵消
模式不同，两棵归约树可能相同或各有优劣。浮点加法不满足结合律，所以改变配对次序足以改变
舍入路径和最终 ULP；算法仍在实数算术意义下计算同一个和。

\editorialnote 英文底本把本题写成“图 A.4 的算法”，但原书图 A.4 是可表示数表，并非
归约算法。结合“避免分支分歧”和第 10 章内容，实际指向原书图 10.7；中文译稿已作同样修正。

\complexity{两种树形归约都是 $O(N)$ 工作、$O(\log N)$ 深度}{$O(N)$ 原地数据或 $O(1)$ 每线程额外状态}{低分歧映射提高 SIMT 利用率，但线程映射并不保证数值稳定性}
\conclusionheading 低控制分歧方案用数值上相距较远的操作数换取连续活跃线程；对升序
非负数据，这一配对通常不利于保留小量贡献。
\discussionheading
\discussionpoint{归约树同时改变关键路径与舍入路径。}
同一批数只是换一棵加法树，为何性能和答案都会变化？顺序累加有 $O(N)$ 的依赖深度，平衡树将其降为 $O(\log N)$，并常把一阶误差增长从约 $O(uN)$ 降到 $O(u\log N)$；这里的 $u$ 是单位舍入误差，但这只是典型上界而非逐输入保证。原书的低分歧映射让更多线程连续活跃，却在首轮把已排序数组的前后半段配对，使小量可能分别消失在大数的 ULP 下。并行数值算法因此不能只优化工作量和占用率，还要把数据到叶节点的映射视为算法语义，并用高精度参考比较不同树的误差分布。

\discussionpoint{稳定求和方法形成精度与成本的谱系。}
如果普通平衡树仍不够准确，还有哪些选择？成对求和保留并行结构；Kahan 或 Neumaier 求和为被舍掉的低位维护补偿量；按指数分桶、超累加器或长定点累加则能把大量项延迟到更精确的表示中合并，甚至支持正确舍入。代价依次可能表现为更多指令、寄存器和共享内存、更复杂的合并状态，或较低的 SIMD/SIMT 利用率，而且补偿更新本身也必须按规定次序实现。数据库聚合、深度学习梯度和科学守恒量对准确度预算不同，合理方案是在给定吞吐、存储和确定性约束下选择谱系中的一点，而非默认最高精度总是最优。

\discussionpoint{排序只有结合符号和条件数才有意义。}
“从小到大相加更准确”是否能无条件套用？对全非负数据，按绝对值从小到大通常让小量先形成可见的部分和；含正负数时，按代数值升序会先聚合大负数，再聚合大正数，可能产生很大的中间量，最后一次抵消仍会暴露输入误差。按绝对值排序也不是万能定理，因为有意让相反符号的近等量先抵消有时反而减少中间范围，效果取决于和的条件数与分布。排序还需额外的 $O(N\log N)$ 工作和数据搬移，所以工程上更常采用分桶或局部重排，并以真实数据、高精度参考和残差共同决定其收益。

\discussionpoint{可复现性与高准确度是两个独立目标。}
同一归约多跑一次得到末位不同，是否就说明某次更不准确？原子更新、动态任务调度或不同分区可能改变合法的加法顺序，从而得到多个都接近参考值的结果；固定树可保证逐位复现，却可能稳定地重复一个误差较大的配对。反过来，补偿或高精度累加能改善误差，但若合并次序仍不固定，结果依旧未必 bitwise 相同。审计、回归定位和某些迭代求解器重视确定性，探索性训练可能更重视吞吐与统计等价；接口应分别声明顺序契约、误差界和性能，而实验记录还应固定分区、工具链及设备拓扑。
\variantheading 在 binary32 中取 $[2^{-24},2^{-24},1]$。若先分别把两个小数加到 1，
每次都落在 tie 并按偶数舍回 1；若先算两个小数，得到 $2^{-23}$，再加 1 得到可表示的
$1+2^{-23}$。两种次序相差 1 ULP，直接展示结合律失效。
\practiceheading 需要跨运行或跨 GPU 可复现时，应固定归约树、块大小和分区顺序，并把
确定性作为接口契约；只要求统计准确时，可在性能预算内选择分桶、Kahan/Neumaier 补偿或
更高精度累加器，并用高精度参考计算 ULP 分布。
\sourcecheck{https://docs.nvidia.com/cuda/floating-point/index.html}{NVIDIA CUDA Floating Point 官方资料对加法非结合性和并行归约差异的说明；对抗检查：排序应按绝对值理解数值稳定性，含正负数据时仅按代数值升序并不能保证相邻项同量级}
\end{solutionexercise}

\begin{solutionexercise}{5}
\problemheading 假设某算术单元以迭代近似算法实现 \code{sin()} 函数，每个时钟周期为
结果再产生 2 个准确的尾数位。架构师决定让该算术函数迭代 9 个时钟周期。假设硬件把剩余
尾数位全部填为 0。对于 IEEE 单精度数，这一 \code{sin()} 硬件实现的最大 ULP 误差是
多少？假设被省略的 ``$1.$'' 尾数位也必须由算术单元产生；亦即 binary32 的有效精度按
隐含首位加 23 个显式 fraction 位共 24 位计。

\approachheading IEEE binary32 的有效数是 24 位：1 个隐含首位和 23 个显式 fraction 位。
九周期只确定 $2\times9=18$ 个最高有效位。

\answerheading 未确定并清零的低位数为
\[
  24-18=6.
\]
硬件结果的相邻量化级相隔 $2^6=64$ 个 binary32 末位单位。因此，以精确实数为参照，
向下清零造成的误差满足
\[
  0\le |\widehat{s}-s|<64\ \mathrm{ULP},
\]
上确界为 64 ULP。若参照“正确舍入到 binary32”的结果，精确值尾部接近下一个粗量化级时，
正确舍入会进位而该硬件仍截断到当前级，二者可相差整整 64 ULP；所以题目通常要求的最大值
是 \textbf{64 ULP}。只观察一个未进位的正确舍入尾数字段会得到 63 ULP，但那漏掉了跨粗
量化边界的进位情形。

该推导假设正常有限结果，并以结果所在 binade 的 binary32 间距定义 ULP。零点附近的
反规格化处理、精确的 $\sin$ 参数归约误差以及 NaN 都不在题设的“唯一误差来自迭代截断”
模型内。

\conclusionheading 18 个准确有效位比 binary32 的 24 位少 6 位，故最大误差规格为
$2^6=64$ ULP（对实数误差是严格小于该值的上确界）。
\discussionheading
\discussionpoint{初等函数的误差来自一整条实现管线。}
把低 6 位清零得到 64 ULP 后，能否把它直接当作任意 \code{sin} 实现的总误差？实际函数通常先判别特殊值，再做周期和象限的参数归约，在小区间使用查表、多项式或迭代近似，随后重构符号并舍入；近似系数、有限精度运算与末端舍入都会增加误差。本题特意假定前 18 位完全准确，因此只隔离了最后截断这一环，不能涵盖前面任何误差。数学库设计常为各阶段分配独立预算，再用组合界或穷举验证总误差；这种分层方法也适用于指数、对数和倒数平方根等硬件函数。如果前级预算未受控，末端多迭代一轮也无法给总误差兜底。

\discussionpoint{巨大输入让参数归约成为主导难题。}
为什么一个看似简单的周期函数会为巨大 $x$ 准备慢速路径？计算 \code{sin(x)} 前必须确定 $x$ 属于哪个 $\pi/2$ 区间并求出很小的余量；当 $x$ 很大时，普通精度的 $x-k\pi/2$ 是两个大数相减，常数低位或整数 $k$ 的一位错误都可能选错象限，其误差远大于末端 64 ULP。高质量实现会保存更多位的 $2/\pi$，以整数式乘法提取象限，并仅在小参数上使用快速近似。快速路径与精确慢速路径的分界需要兼顾延迟、代码体积和最坏误差，也说明输入范围是函数准确度契约不可省略的条件。

\discussionpoint{数学库必须声明准确度契约。}
用户看到同名的 \code{sin}，能否假定不同编译选项都给出相同末位？正确舍入要求结果等于精确实数按目标模式舍入后的唯一浮点数；faithful rounding 通常只保证落在包围精确值的相邻候选之一，而快速数学路径可能仅承诺更宽的 ULP 或相对误差，并放宽 NaN、反规格化或特殊值行为。前者需要昂贵的难例处理，后者可换取吞吐和更短延迟。金融定价、可证明科学模拟与机器学习推理承担的风险不同，API、编译标志和测试报告都应明确契约，调用方再决定是否以性能交换准确度。

\discussionpoint{最坏误差验证不能依赖随机平均。}
随机输入的平均 ULP 很小，是否足以证明 64 ULP 上界没有被突破？高误差点可能集中在零点、极值、binade 边界、舍入中点和参数归约换路处，随机采样命中这些狭窄集合的概率极低；应使用位型枚举、邻接浮点数及多精度参考主动搜索，并覆盖反规格化与非有限值。对“只截取精确展开的前 24 位”这一变体，必须构造正确舍入向上进位的例子，才能验证它与正确舍入结果至多相差 1 ULP、且该界确实可达到。最后还要分别测量指令延迟和流水吞吐，因为九周期延迟并不意味着执行单元每九周期才接受一个新操作。
\variantheading 若迭代 10 个周期，则得到 20 个准确有效位，还缺 $24-20=4$ 位，最大
规格变为 $2^4=16$ ULP；若迭代 12 个周期得到的是已经正确舍入的 24 位 binary32 结果，
则末端清零误差为 0。若硬件只是截取精确展开的前 24 位而没有 guard/sticky 位，仍可能与
正确舍入结果相差至多 1 ULP，并且在正确舍入选择相邻上端点时可以恰差 1 ULP；“得到 24 位”
本身不足以推出零 ULP。
\practiceheading 选择快速数学函数时，应在业务输入域上同时测延迟、吞吐量和最大/分位 ULP，
并专门覆盖零点、极值、巨大参数与 NaN/无穷。编译选项和 GPU 架构会改变实现路径，验证结果
应与工具链及目标架构一起归档。
\sourcecheck{https://standards.ieee.org/ieee/754/6210/}{IEEE 754-2019 官方标准页面；对抗检查：若忘记隐含首位会误算为少 5 位并得到 32 ULP，若把截断尾部最大位型机械写成 63 又会漏掉正确舍入的进位}
\end{solutionexercise}
